%-------------------------------------------------------------------------------
%	ABSTRACT PAGE
%-------------------------------------------------------------------------------

\addchaptertocentry{Хураангуй} % Хураангуйг гарчигт нэмэх

\begin{center}
{\scshape\Large \univname\par} % Их сургуулийн нэр
{\scshape\large \facname\par}\vspace{0.5cm} % Их сургуулийн нэр
{\huge\textbf{{Хураангуй}} \par}
\bigskip
{\Large{\ttitle} \par} % Тезисийн нэр
\bigskip

{\normalsize \shortname \par}\addressname

% Зохиогчийн нэр

\end{center}

\textit{\textbf{Түлхүүр үгс: \keywordnames}}
\paragraph{}Интернэтийн хэрэглэгчдийн тоо дэлхий даяар хурдацтай нэмэгдэж буй өнөө үед цахим орчинд тархаж буй хортой агуулга, кибер дарамт, үзэн ядалтын хэл зэрэг сөрөг үзэгдлүүд нийгмийн ноцтой асуудал болж байна. Гэвч сөрөг агуулга нь бүгд хориглох шаардлагатай биш --- бүтээлч шүүмжлэл, асуудалд суурилсан шинжилгээ нь ардчилсан нийгмийн чухал суурь болдог. Иймд автомат шүүлтийн систем нь хортой агуулга (Toxic Negative) болон бүтээлч шүүмжлэл (Constructive Criticism)-ийг нарийн ялгаж чаддаг байх шаардлагатай.

Энэхүү судалгаанд Монгол хэл дээрх цахим орчны агуулгыг хоёр шатлалт (two-stage) NLP системээр ангилах аргыг боловсруулсан. MN-BERT загварын fine-tuning, Монгол хэлний agglutinative морфологид тохирсон preprocessing pipeline, болон Gradio-д суурилсан хэрэглэгчийн интерфейсийг нэгтгэсэн. Дөрвөн ангилалт (Positive / Neutral / Toxic Negative / Constructive Criticism) шошгологдсон 10,000 гаруй дээж бүхий Монгол хэлний сошиал медиа датасет бүрдүүлж, MN-BERT загварыг суурь загваруудтай (Logistic Regression, Naive Bayes, SVM) харьцуулан үнэлсэн.

Туршилтын үр дүнгээр MN-BERT загвар нь нэгдсэн систем дээр $81.43\%$ нарийвчлал (accuracy) болон $0.776$ Macro-F1 үзүүлэлтэд хүрсэн нь TF-IDF суурьтай baseline загваруудаас дунджаар $16\%$ илүү гүйцэтгэл юм. Ялангуяа Toxic болон Constructive ангиллыг ялгах Stage~2-ын нарийвчлал $91.44\%$-д хүрсэн нь системийн практик ач холбогдлыг баталж байна. Энэхүү ажил нь Монгол хэлний NLP-ийн салбарт хортой агуулгыг шүүхдээ иргэдийн үзэл бодлоо илэрхийлэх эрх чөлөөг (бүтээлч шүүмжлэл) хамгаалах техникийн шийдлийг санал болгож буйгаараа шинжлэх ухааны шинэлэг талтай юм.
